% SHARD-ID: RC-08I-RIEMANN-TOMOGRAPHY
% SHARD-TITLE: The Riemann Tomography Problem, Round 1: the Dilation Channel, the Prime-Content Channel, the Kinematic Commutant, and the Graph Control
% SHARD-SUMMARY: The Weil functional is known exactly from primes, pole and archimedean place; treating its positivity as the unknown state of a tomography problem, the round measures what each family of arithmetic observations teaches: along the dilation window the information per Fourier mode saturates at N of order x log x and the e^{-4 pi x} floor is a cancellation of every prime power below x (pole plus archimedean alone is indefinite); along prime content the inter-place structure is of the order of the bump width and is carried by the pole term first.
% SHARD-SUMMARY: The maximum-entropy ansatz in the Loewner coordinates is empty (the admissible next datum is an interval whose unstructured centre is not a Loewner column, and the positive set of prime data with pole and archimedean fixed is unbounded), while in the lattice coordinates positivity alone pins the von Mangoldt weights at the known positions log n to 1e-7; a certified Ramanujan graph control learns along a power law and collapses at a finite window, unlike zeta.
% SHARD-SUMMARY: Sources fixed: Yoshida's unconditional positivity up to support width log 2 is a pole-plus-archimedean statement (the archimedean form alone is indefinite), Landau's formula is what a single-place restriction sees, Burg's extension is the disc centre, and RH is a Pi^0_1 sentence.
% SHARD-KEYWORDS: tomography, Weil form, explicit formula, window, Loewner, Schur complement, mutual information, maximum entropy, Burg, prime content, lattice, Kronecker sum, product state, Schmidt defect, Landau formula, Yoshida, Bombieri, short support, commutant, recession cone, von Mangoldt, Ramanujan graph, calibration, Pi^0_1
% SHARD-NOTES: none
% SHARD-SCRIPTS: scripts/rtp1_prime_content.py, scripts/rtp1_commutant.py, scripts/rtp1_calibration.py

\section{The Riemann Tomography Problem, Round 1}
\label{sec:riemann-tomography}

\emph{Question (TJO, 2026-09-24).} Treat the metric as an unknown, unnormalised quantum state. RH is not
assumed, the zeros are not used, and only arithmetic data are admitted as observations. Which observations
affect the state at all, and how fast does each family contract the prior? Records:
\texttt{notes/metric-tomography/metric-as-state.md}, \texttt{notes/rtp-round-1/} (brief, lanes A1, A2, B1,
B2, V, review). Conventions: the Weil form of Connes--Consani--Moscovici as implemented in \texttt{zst/},
$Q(f,g) = \Psi(f^**g)$, $\Psi = W_{0,2}-W_{\mathbb R}-\sum_pW_p$, on the window $[\lambda^{-1},\lambda]$,
$x = \lambda^2$, $L = \log x$, Fourier basis $V_n$, $|n|\le N$, Loewner matrix $T_{nm} = (b_n-b_m)/(n-m)$,
$T_{nn} = a_n$, with $(a_n,b_n)$ certified balls (pole, prime and archimedean parts in closed form). Zeros of
$\zeta$ enter only the labelled comparison steps.

\subsection{The problem and the two coordinate systems}
An \emph{observation} is a finite family $f_1,\dots,f_n$ of admissible test functions with Gram matrix $G_{ij} = Q(f_i,f_j)$, computable from primes. The \emph{posterior} after an observation is the set of positive
extensions of the data; the \emph{MaxEnt ansatz} is the maximum-determinant extension; the \emph{learning
rate} of a channel is the contraction of the posterior as arithmetic data are added. Two coordinate systems
appear: the Loewner coordinates $(a_n,b_n)$ of the window, and the lattice coordinates (weights at the
positions $\log n$). The round's central structural finding is that they are not equivalent for the
tomography (\cref{obs:rtp-round-1-reading}).

\begin{lemma}[Bordering, mutual information, and the Loewner interval]
\label{lem:bordering-interval}
\claimstatus{lem:bordering-interval}{sketched}
Let $G_K$ be positive definite Hermitian and border it by a column $c$ and diagonal $d$. (i) The bordered
matrix is positive semidefinite iff $d-c^*G_K^{-1}c\ge0$; with $d$ known, the maximum-determinant completion is
$c = 0$, and the log-determinant gain of the true column over it, $\varDelta I = \log\big(d/(d-c^*G_K^{-1}c)\big)\ge0$,
is twice the Gaussian mutual information between the new and the old modes. (ii) In the window form the
bordering $|n|\le N\to N+1$ adds one even and one odd basis vector whose columns and diagonals are affine in the
single new datum $b_{N+1}$ given $a_{N+1}$; each Schur complement is a concave quadratic in
$b_{N+1}$, the admissible $b_{N+1}$ of each block is an interval, its centre is the maximum-determinant value,
the true value lies on the boundary iff the bordered block is singular, and the joint admissible set is the
intersection of the two intervals. (iii) The unconstrained point $c = 0$ is not a Loewner column unless
$b_1=\dots=b_N = 0$.
\end{lemma}

Part (ii) is the analogue of \cref{prop:extension-disc}(i)--(iv) for the non-Toeplitz window; (iii) says that
the naive MaxEnt prediction of the tomography note is structurally excluded. Proof: lane A1, section A1.1.

\begin{numerical}[The dilation channel: resolution saturates at $N\sim1.7\,x\log x$; the floor is a cancellation]
\label{num:rtp-dilation-channel}
\claimstatus{num:rtp-dilation-channel}{numerical}
Certified in ball arithmetic (\texttt{zst/tools/rtp1\_a1.c}, outputs \texttt{outputs/rtp1\_a1\_*.txt}).
(a) At $x = 13, 25, 50$ the log-determinant of the even block is minimal at $N = 56, 134, 352$
($N_{\mathrm{sat}}\approx1.7\,x\log x$; the benchmark's $7.5x$ holds only near $x = 50$). Below $N_{\mathrm{sat}}$
each new mode pair carries tens of nats and the true $b_{N+1}$ sits at the edge of its interval; above, a
fraction of a nat, interior, with interval half-width $0.5$ to $3.7$ at every $x$. (b) The minimal even
eigenvalue at saturated $N$ is $3.48\cdot10^{-59}$, $2.43\cdot10^{-123}$, $2.27\cdot10^{-257}$, falling
$5.34$--$5.36$ digits per unit of $x$ against $5.46$ for $e^{-4\pi x}$. (c) Comparison step (zeros used): the
certified first-zero error is $7.0\cdot10^3$, $9.3\cdot10^3$, $1.08\cdot10^4$ times the minimal eigenvalue at
$x = 13, 25, 50$ and has the same slope in $x$ at every fixed $N$. (d) At fixed window $L = \log50$, $N = 60$,
pole plus archimedean alone has three negative even eigenvalues; adding prime powers one at a time leaves
$3$ to $16$ negative eigenvalues until the last one, $49$, enters (with all but $49$: minimal eigenvalue
$-5.66\cdot10^{-7}$; with it, $1.75\cdot10^{-116}$). Each prime power below $x = 50$ contributes at least
$10^{34}$ times the final minimal eigenvalue: the floor is a cancellation residual of all of them.
\end{numerical}

\subsection{The prime-content channel}

For a finite prime set $S$ and exponent box $\{0\le a_p\le A\}$, take a $C_c^\infty$ bump of half-width
$\delta$ in the log coordinate centred at each $\log\prod_pp^{a_p}$. Admissibility: $\delta$ small enough that
$f^**g$ touches no prime power other than the lattice ratios themselves; the largest admissible $\delta$ is
half the distance from the lattice ratios to the nearest foreign prime power (binding pairs $6/7$, $36/37$,
$108/109$, $1296/1297$ for $\{2,3\}$, $A = 1,\dots,4$; $30/31$, $450/449$ for $\{2,3,5\}$, $A = 1,2$). A prime
term enters a Gram entry only when the two lattice points differ in at most one coordinate; every
\emph{mixed} entry is pole plus archimedean.

\begin{lemma}[Zeroing the mixed entries gives the Kronecker sum]
\label{lem:mixed-entries-kronecker}
\claimstatus{lem:mixed-entries-kronecker}{sketched}
With admissible $\delta$, replacing every Gram entry between lattice points differing in two or more
coordinates by zero yields exactly the Kronecker sum of the one-prime forms (the shared diagonal counted
once); its minimal eigenvector is the product of the one-prime minimal eigenvectors.
\end{lemma}

\begin{numerical}[The prime-content channel is perturbative in the bump width]
\label{num:rtp-prime-content}
\claimstatus{num:rtp-prime-content}{numerical}
Certified balls at 200 bits (\texttt{scripts/rtp1\_prime\_content.py}, 192 checks; two diagonal archimedean
constants by quadrature with error $<3\cdot10^{-55}$, a multiple of the identity). The general implementation
reproduces \texttt{zst}'s $(a_n,b_n)$ to $2.2\cdot10^{-60}$. All 42 Gram matrices ($S = \{2\},\{2,3\},\{2,3,5\}$,
$A\le6,4,4$, $\delta = 0.9,0.3,0.1$ of the maximum) are positive definite with simple minimal eigenvalue. The
mixed entries raise the minimal eigenvalue by an amount linear in $\delta$; the Schmidt defect of the minimal
eigenvector across $\{p\}\,|\,$rest and the eigenvector movement on adding a prime ($\{2\}\to\{2,3\}$:
$4.6\cdot10^{-5}$ to $7.6\cdot10^{-10}$) scale as $\delta^2$ (fitted slopes $2.00$--$2.06$). Inside the mixed
entries the pole part exceeds the archimedean part by a factor $2.4$ to $21$, with opposite sign; removing the
pole part alone makes $\{2,3\}$, $A = 1$ indefinite. Comparison step (zeros used): the prime-side values agree
with the zero-side sums over 2000 certified zeros to $3\cdot10^{-18}$ at $\delta = 0.1$, with truncation error
$\sim\exp(-2.2\sqrt{\delta\gamma_M})$.
\end{numerical}

\begin{observation}[What a single-place restriction sees]
\label{obs:single-place-restriction}
\claimstatus{obs:single-place-restriction}{sketched}
\cref{obs:prime-by-prime-blind} concerns single-place \emph{ansätze}. The \emph{restriction} of the true
form to test functions $g = \sum_kc_k\,\phi(x-k\log p)$ has $\hat g(s) = \hat\phi(s)P(p^{-s})$ with $P$ a
polynomial, so $Q(g,g) = \sum_\zero|\hat\phi(\zero)|^2|P(p^{-\zero})|^2$: it sees the zeros folded on the
$p$-circle, the content of Landau's formula (\cref{cit:landau-formula}); an off-line pair is visible to it
unless $P$ vanishes there. Mixed $\{p,q\}$ entries carry no prime term and their zero-side Landau main term
vanishes, so the two-prime channel measures the error term of Landau's formula at composite $x$; by
\cref{num:rtp-prime-content} that error term enters the eigenvector only at order $\delta^2$.
\end{observation}

\subsection{Sources: short support, Landau, Burg, and the $\Pi^0_1$ form}

\begin{citedfact}[Yoshida: unconditional positivity up to support width $\log2$]
\label{cit:yoshida-short-support}
\claimstatus{cit:yoshida-short-support}{cited}
For $a = \tfrac12\log2$ the Weil form is nonnegative on every smooth test function supported in $[-a,a]$ in
the log coordinate, with equality only for the zero function; the proof is computer-assisted. On that space
the prime sum vanishes, so this is a statement about the pole-plus-archimedean form.


\par\noindent Quoted (\texttt{prov:yosh92-thm1}):
{\small\texttt{\detokenize{Theorem 1. Let a= log 2/2. We have}}}
\end{citedfact}

\begin{citedfact}[Bombieri: a lower bound for short support]
\label{cit:bombieri-short-support}
\claimstatus{cit:bombieri-short-support}{cited}
For $F$ supported in an interval $I$ of length $|I|<\log2$, the Weil form satisfies
$\ge\big(\log(1/|I|)-\log^+\log(1/|I|)-O(1)\big)\|F\|^2$, unconditionally, with the $O(1)$ unspecified.


\par\noindent Quoted (\texttt{prov:bomb00-thm12}):
{\small\texttt{\detokenize{Theorem 12. If F (x) has compact support in an interval I of length |I | < log 2 we have}}}
\end{citedfact}

\begin{citedfact}[The archimedean form alone is indefinite]
\label{cit:archimedean-symbol-sign}
\claimstatus{cit:archimedean-symbol-sign}{cited}
The Fourier symbol of the archimedean term, $-\log\pi+\mathrm{Re}\,(\Gamma'/\Gamma)(\tfrac14+\tfrac{it}2)$, is
negative for $|t|<6.2898$; on the even box of width $\log2$ the archimedean part is $-1.217$ and the pole part
$+1.400$. Positivity on short support is a property of the two together.


\par\noindent Quoted (\texttt{prov:cc20-hplus}):
{\small\texttt{\detokenize{h_+(\tau)=-\log \pi + \Re(\lambda(1/4 +i\tau/2)), \qquad \lambda(z)=\Gamma'(z)/\Gamma(z).}}}
\end{citedfact}

\begin{citedfact}[Landau's formula]
\label{cit:landau-formula}
\claimstatus{cit:landau-formula}{cited}
For fixed $x>1$, $\sum_{0<\gamma\le T}x^{\zero} = -\tfrac{T}{2\pi}\vM(x)+O(\log T)$; the main term is present
iff $x$ is a prime power. Gonek's version is uniform in $x$ and $T$.


\par\noindent Quoted (\texttt{prov:landau12-satz1}):
{\small\texttt{\detokenize{Satz 1: Es ist bei festem «>1}}}
\end{citedfact}

\begin{citedfact}[Burg: the disc centre is the maximum-entropy extension]
\label{cit:burg-maxent}
\claimstatus{cit:burg-maxent}{cited}
Given a positive Toeplitz window, the admissible next autocorrelation fills a disc; choosing its centre, i.e.
zero reflection coefficients from then on, yields the maximum-entropy spectrum, which exists and is unique.


\par\noindent Quoted (\texttt{prov:burg75-centre}):
{\small\texttt{\detokenize{if one extends the autocorrelation function to infinity by always selecting the center value for each consecutive lag value, one generates the autocorrelation function corresponding to the maximum entropy spectrum.}}}
\end{citedfact}

\begin{citedfact}[Maximum-determinant completion]
\label{cit:maxdet-completion}
\claimstatus{cit:maxdet-completion}{cited}
A partial Hermitian matrix with a positive definite completion has a unique maximum-determinant completion,
characterised by its inverse vanishing on the unspecified entries (Dempster; Grone--Johnson--S\'a--Wolkowicz).


\par\noindent Quoted (\texttt{prov:va15-maxent-completion}):
{\small\texttt{\detokenize{The solution W is also called the maximum entropy completion of A, since it maximizes the entropy}}}
\end{citedfact}

\begin{citedfact}[RH is a $\Pi^0_1$ sentence]
\label{cit:rh-pi01}
\claimstatus{cit:rh-pi01}{cited}
RH holds iff $\sum_{d\mid n}d\le H_n+e^{H_n}\log H_n$ for every $n\ge1$ (Lagarias), and iff an explicit Diophantine
inequality holds for every $n$ (Davis--Matiyasevich--Robinson); there is an explicit Turing machine that halts
iff RH is false.


\par\noindent Quoted (\texttt{prov:lag02-thm}):
{\small\texttt{\detokenize{Problem $E$ is equivalent to the Riemann hypothesis.}}}
\end{citedfact}

\begin{citedfact}[Gonek: Landau's formula uniform in $x$ and $T$]
\label{cit:gonek-uniform-landau}
\claimstatus{cit:gonek-uniform-landau}{cited}
Landau's formula holds with an error term uniform in $x$ and $T$, growing linearly in $x$.


\par\noindent Quoted (\texttt{prov:klp00-gonek}):
{\small\texttt{\detokenize{this direction was obtained by Gonek [4], [5]. He proved that for T > 1}}}
\end{citedfact}

\begin{citedfact}[Bombieri: negative eigenvalues of the Weil functional on longer intervals]
\label{cit:bombieri-negative-eigenvalues}
\claimstatus{cit:bombieri-negative-eigenvalues}{cited}
The quadratic form has negative eigenvalues once the support is long enough; positivity is a small-support
statement only.


\par\noindent Quoted (\texttt{prov:bomb00-neg-eigen}):
{\small\texttt{\detokenize{if the Riemann Hypothesis is false but only with finitely many non-trivial zeros off the critical line we show that the number of negative eigenvalues is precisely one-half of the number of zeros failing to satisfy the Riemann Hypothesis, provided the truncation is big enough.}}}
\end{citedfact}

\begin{citedfact}[The Li coefficients in closed form]
\label{cit:li-arithmetic-formula}
\claimstatus{cit:li-arithmetic-formula}{cited}
$\lambda_n = -\sum_m\binom nm\eta_{m-1}+\sum_m(-1)^m\binom nm(1-2^{-m})\zeta(m)+1-\tfrac n2(\gamma+\log\pi+2\log2)$,
with $\eta_j$ the Laurent coefficients of $\zeta'/\zeta$ at $s = 1$ (Bombieri--Lagarias, via Coffey).


\par\noindent Quoted (\texttt{prov:coffey05-li-formula}):
{\small\texttt{\detokenize{$$\lambda_n=-\sum_{m=1}^n {n \choose m} \eta_{m-1}+\sum_{m=2}^n (-1)^m {n \choose m} (1-2^{-m})\zeta(m)+1-{n \over 2}(\gamma +\ln \pi + 2\ln 2), \eqno(10)$$}}}
\end{citedfact}

\begin{citedfact}[Hulth\'en: the Heisenberg chain energy]
\label{cit:hulthen-energy}
\claimstatus{cit:hulthen-energy}{cited}
The ground-state energy of the spin-$\tfrac12$ antiferromagnetic Heisenberg chain $H = \sum S_i\cdot S_{i+1}$
is $N(\tfrac14-\log2)$ in the thermodynamic limit.


\par\noindent Quoted (\texttt{prov:franchini17-hulthen}):
{\small\texttt{\detokenize{This result was originally derived by Hulthen \cite{hulthen38}.}}}
\end{citedfact}

\begin{citedfact}[Tarski--Seidenberg]
\label{cit:tarski-seidenberg}
\claimstatus{cit:tarski-seidenberg}{cited}
Projections of semialgebraic sets are semialgebraic, and first-order definable sets over a real closed
field are semialgebraic.


\par\noindent Quoted (\texttt{prov:coste02-ts-second}):
{\small\texttt{\detokenize{Theorem 2.3 (Tarski-Seidenberg – second form) Let A be a semialge- braic subset of Rn+1 and π : Rn+1 → Rn , the projection on the first n coordi- nates. Then π(A) is a semialgebraic subset of Rn .}}}
\end{citedfact}

\begin{proposition}[Burg's extension is the maximum-determinant completion of the next lag]
\label{prop:burg-is-maxdet}
\claimstatus{prop:burg-is-maxdet}{sketched}
For Toeplitz data $\nuseq{0},\dots,\nuseq{K}$ with $\Wform_K$ positive definite, the maximum-determinant
positive completion of the next lag is the centre $c_K$ of the disc of \cref{prop:extension-disc}, and
coincides with Burg's zero-reflection-coefficient extension.
\end{proposition}

\subsection{The kinematic commutant and the two coordinate systems}

\begin{proposition}[Dimension of the kinematic space; positivity removes nothing]
\label{prop:window-kinematic-dimension}
\claimstatus{prop:window-kinematic-dimension}{sketched}
In the basis $V_n$, $|n|\le N$: Hermitian forms commuting with the reflection $V_j\mapsto V_{-j}$ form a space
of dimension $(N+1)^2+N^2$; real ones $(N+1)(2N+1)$; Loewner forms with the fixed vector $\eta = \sum V_n$ have
dimension $4N+1$ and are automatically real; Loewner and reflection-commuting: $2N+1$, exactly the data
$(a_0,\dots,a_N,b_1,\dots,b_N)$. The identity is interior to the positive cone of the last space, so
positivity removes no dimension before arithmetic data enter.
\end{proposition}

\begin{numerical}[MaxEnt is empty in Loewner coordinates; positivity pins the weights in lattice coordinates]
\label{num:rtp-commutant}
\claimstatus{num:rtp-commutant}{numerical}
$x = 13$, $N = 20$ (\texttt{scripts/rtp1\_commutant.py}, 37 checks; mpmath at 100 digits for positions,
double-precision conic programmes with duality gaps $\sim4\cdot10^{-9}$ re-verified at 100 digits). With pole
and archimedean data fixed and the prime data $(u,v)$ free: the positive set is full-dimensional and
unbounded (recession cone the whole positive cone), so no maximum-determinant element exists; zero prime data
are inadmissible (minimal eigenvalue $-1.96$); the minimal-norm admissible element has norm $2.36$ against the
truth's $8.21$ and captures $8.3\%$ of the true prime data ($2.8\%$ at $N = 60$); the prime-number-theorem
mean captures $22\%$ and is inadmissible. The truth sits on the boundary ($1.6\cdot10^{-39}$), isolated on
its ray (admissible multiples in $[1-5.7\cdot10^{-36},\,1+2.0\cdot10^{-38}]$). With the positions $\log n$,
$n = 2,\dots,12$, fixed and only the weights free, positivity confines each weight to a box of width
$\le1.1\cdot10^{-7}$ for $n\le9$ ($1.3\cdot10^{-3}$ at $n = 12$), including the zero weights at $6, 10, 12$.
\end{numerical}

\begin{numerical}[The graph control: power law and collapse]
\label{num:rtp-calibration}
\claimstatus{num:rtp-calibration}{numerical}
A random cubic graph on 40 vertices (seed 20260924), certified Ramanujan by exact arithmetic, $R = 78$ atoms
(\texttt{scripts/rtp1\_calibration.py}, 33 checks; certified eigenvalue enclosures at 600 bits). Along the
window $K$ (also the data cutoff), the minimal eigenvalue falls as a power law in $x_{\mathrm{eq}} = 2^K$
($0.13$--$0.79$ digits per unit $K$) and reaches exactly zero at $K = R$; the next trace is interior to its
disc for $K<77$ and on the boundary at $K = 77$; at fixed window every partial form is indefinite until the
last cycle length enters; the minimal eigenvalue is a cancellation residual there too; at the largest window
the positive set of prime data is bounded and its maximum-determinant element is the graph prime number
theorem to $4.7\cdot10^{-10}$. Comparison step: the kernel-root error tracks the minimal eigenvalue as its
$0.88$ power.
\end{numerical}

\begin{observation}[The round's reading]
\label{obs:rtp-round-1-reading}
\claimstatus{obs:rtp-round-1-reading}{sketched}
(i) The prime-content channel has no contracting signal at matched information; its inter-place effects
are $O(\delta)$ in eigenvalues and $O(\delta^2)$ in eigenvectors, with admissibility forcing $\delta$ down like
the inverse of the largest lattice ratio. (ii) The off-axis weight of the two-prime eigenvector is product-state
weight; the genuine correlation is the $O(\delta^2)$ Schmidt defect, carried by pole plus archimedean, pole
first. (iii) The $e^{-4\pi x}$ law is arithmetic in the sense of the round: the kinematics fix
$N_{\mathrm{sat}}$ (a zero-counting scale), pole plus archimedean alone is indefinite, the floor is a
cancellation of every prime power below $x$, and the graph control shows a power law and a collapse instead;
whether the rate is a property of ``a spectrum with $\zeta$'s Weyl law'' or of these zeros needs a growing
family of controls. (iv) The maximum-entropy ansatz of the tomography note is empty in the Loewner
coordinates (\cref{lem:bordering-interval}(iii), \cref{num:rtp-commutant}); the informative coordinates are
the weights at the known positions $\log n$, where positivity alone pins the primes. The tomography problem
is therefore posed in the lattice coordinates, with the window as the resolution and $\delta$ as the price of
prime purity.
\end{observation}

\paragraph{Not established, and what is next.} No statement here bears on RH. The lattice-coordinate pinning
is a rigidity statement about a nearly singular form; how its box widths scale with $N$ and $x$ is the next
measurement, and the Landau error term of \cref{obs:single-place-restriction} is what a wider, impure
prime-content window would measure. Sources are byte-cited in \texttt{db/provenance.tsv}.
